%!TEX encoding = UTF8
%!TEX root = 0-notes.tex

\chapter{Fonctions réelles}
\label{chap:3}


\section{Comparaisons asymptotiques}

\subsection{Petit o}

\subsection{Grand O}

\subsection{Équivalence}

\section{Développements limités}

\subsection{Premier ordre}

\subsection{Théorèmes de Taylor-Lagrange et Taylor-Young}

\subsection{Applications}

\section{Régularité}

\subsection{Continuité et continuité séquentielle}

\subsection{Théorème des valeurs extrêmes}

\subsection{Continuité uniforme et théorème de Heine-Borel}

\subsection{Fonctions $k$-lipschitziennes et applications}

\section*{Exercices supplémentaires}
\addcontentsline{toc}{section}{Exercices supplémentaires}

\exe{,difficulty=0}{
	Montrer le développement limité
		\[ ln(1-x) = - \left(\frac1x + \frac2{x^2} + \frac3{x^3} + \ldots + \frac{n}{x^{n}} \right) + O\left(\frac1{x^{n+1}}\right). \]
	% suite exo série 1/n = ln(n) + gamma + O(1/n) ?
}{exe:DL-ln}{
	todo
}

%\section*{Solutions}
%\addcontentsline{toc}{section}{Solutions}
%
%\shipoutAnswer



